% Compile: tectonic -X compile v1.0.1.tex
\documentclass[10pt,a4paper]{article}

\usepackage[left=20mm,right=20mm,top=18mm,bottom=18mm]{geometry}
\usepackage{fontspec}
\setmainfont{TeXGyreTermesX}[Extension=.otf, UprightFont=*-Regular, BoldFont=*-Bold, ItalicFont=*-Italic,
  BoldItalicFont=*-BoldItalic]
\usepackage{amssymb}
\usepackage{graphicx}
\usepackage{array}
\usepackage{enumitem}
\usepackage{titlesec}
\usepackage{float}
\usepackage{caption}
\usepackage[hidelinks]{hyperref}

\setlength{\parindent}{0pt}
\setlength{\parskip}{4pt}
\pagestyle{empty}
\titleformat{\section}{\normalsize\bfseries\scshape}{\Roman{section}.}{0.5em}{}
\titlespacing*{\section}{0pt}{9pt}{3pt}
\setlist[itemize]{leftmargin=1.2em, itemsep=1pt, parsep=0pt, topsep=1pt}
% IEEE style: "TABLE I" above in small caps, "Fig. 1." below
\renewcommand{\thetable}{\Roman{table}}
\captionsetup[table]{position=top, labelsep=newline, font={footnotesize,sc}, justification=centering, skip=3pt}
\captionsetup[figure]{labelsep=period, font=footnotesize, name=Fig., skip=2pt}
\renewcommand{\arraystretch}{1.15}



\begin{document}

\begin{center}
  {\LARGE\bfseries Release Notes}\par\vspace{3pt}
  {\small 30 October 2026}
\end{center}

\section{ROS support}
Version v1.0.1 adds optional ROS 1 and ROS 2 nodes for image and IMU topics. Both nodes use the same OpenVIO estimator
and the same input processing code. Monocular and stereo configurations are available for EuRoC MAV and TUM-VI.

\begin{table}[H]
\centering
\caption{ROS distributions}
\begin{tabular}{lll}
\hline\hline
ROS & Distribution & Ubuntu \\
\hline
ROS 1 & Noetic & 20.04 \\
ROS 2 & Humble & 22.04 \\
ROS 2 & Jazzy & 24.04 \\
\hline\hline
\end{tabular}
\end{table}

\section{Separate builds}
Use \texttt{./build\_ros1.sh} for ROS 1. Use \texttt{./build\_ros2.sh humble} or
\texttt{./build\_ros2.sh jazzy} for ROS 2. The two scripts run their respective build tools directly.
Each distribution has its own output under \texttt{build\_ros/<distro>/}. Source
\texttt{build\_ros/<distro>/install/setup.bash} in each terminal before use.

\begin{itemize}
\item The standalone build remains \texttt{./build.sh}, with output in \texttt{build/}.
\item The ROS build does not call the standalone script or write to its output directory.
\item Noetic uses a local Eigen 3.4 installation. The system Eigen installation stays unchanged.
\item The ROS build stops if a dependency install, configuration step or compiler command fails.
\end{itemize}

\section{Run and display}
Run one command for the installed ROS distribution:
\begin{quote}\small\ttfamily
roslaunch openvio\_ros1 openvio.launch dataset:=euroc\\
ros2 launch openvio\_ros2 openvio.launch.py dataset:=euroc
\end{quote}
Start the node before playback. For ROS 1 bags, use \texttt{rosbag play -d 2 <bag>}.
Convert ROS 1 bags before ROS 2 playback.

The launch files open RViz with the estimated path, landmarks, tracked image and the \texttt{world} to \texttt{imu}
transform. Use \texttt{stereo:=true} for a stereo dataset or \texttt{config:=/path/to/sensor.yaml} for another sensor.
The optional \texttt{output\_path} argument writes the IMU trajectory in TUM format.

\section{Verification and scope}
ROS 1 Noetic and ROS 2 Jazzy were built and checked with EuRoC MH\_01\_easy monocular data. The checks confirmed a
growing path, changing feature images, nonempty landmarks, and coordinate transforms. RViz displayed the outputs.
These checks verify the ROS interface and display. They do not provide a new accuracy benchmark.

ROS 2 Humble is not tested. Stereo and live sensor input were not rerun in this release check. The accuracy results
in the v1.0.0 release notes remain the reference results.

\section{Documentation}
The documentation links now follow this order: Quick start, ROS 1 / ROS 2, Release notes. See
\href{https://github.com/93won/openvio/blob/master/doc/ros.md}{\texttt{doc/ros.md}} for build commands, topic names,
launch arguments, RViz controls and fault checks.
\end{document}
